\section*{Overview}

The greatest common divisor is one of the smallest ideas in number theory, but it sits underneath modular inverses,
CRT, fraction normalization, and linear Diophantine equations. Extended Euclid is the version worth remembering in
contests because it produces coefficients, not just the gcd itself.

\section*{When to Use It}

Use this when:

\begin{itemize}
  \item divisibility is part of the condition,
  \item you need to simplify a fraction or a ratio,
  \item the equation looks like \(ax + by = c\),
  \item a modular inverse exists only when the gcd condition is satisfied.
\end{itemize}

\section*{Core Idea}

Euclid's algorithm repeatedly replaces \((a, b)\) by \((b, a \bmod b)\). The gcd does not change under that step, and
the numbers get smaller quickly.

Extended Euclid keeps track of coefficients \(x, y\) such that
\[
ax + by = \gcd(a, b).
\]
Those coefficients are exactly what later number-theory algorithms need.

\section*{Key Insight}

The recursive remainder relation
\[
\gcd(a, b) = \gcd(b, a \bmod b)
\]
is not just a gcd trick. It lets every remainder step be unwound into a linear combination of the original numbers.
That is why extended Euclid can recover the coefficients.

\section*{Operations / Main Technique}

\begin{itemize}
  \item \textbf{gcd}: decide divisibility and reduce ratios.
  \item \textbf{lcm}: compute by \(a / \gcd(a,b) \cdot b\) to avoid overflow as much as possible.
  \item \textbf{extended gcd}: recover coefficients.
  \item \textbf{Diophantine existence}: \(ax + by = c\) is solvable iff \(\gcd(a,b)\) divides \(c\).
\end{itemize}

\paragraph{Worked example.}
If \(\gcd(18, 30) = 6\), then \(18x + 30y = 12\) is solvable, but \(18x + 30y = 5\) is not.

\section*{Correctness Intuition}

Each Euclid step preserves the set of common divisors, so it preserves the gcd. Extended Euclid simply records how the
current pair came from the previous pair. Unwinding those steps reconstructs the linear combination for the original
inputs.

\section*{Complexity Analysis}

Euclid and extended Euclid run in \(O(\log \min(a, b))\) time and \(O(\log \min(a, b))\) recursion depth.

\section*{Implementation}

The code includes:

\begin{itemize}
  \item \texttt{gcd\_ll},
  \item \texttt{lcm\_ll},
  \item \texttt{ext\_gcd} producing coefficients,
  \item a small solver for \(ax + by = c\).
\end{itemize}

\section*{Common Pitfalls}

\begin{itemize}
  \item Forgetting that a modular inverse requires gcd \(= 1\).
  \item Overflow in the direct formula \(a \cdot b / \gcd(a,b)\).
  \item Losing sign conventions when negative numbers are allowed.
  \item Treating one particular Diophantine solution as the only one.
\end{itemize}

\section*{Variants / Extensions}

\begin{itemize}
  \item Modular inverse via extended Euclid.
  \item Chinese remainder theorem.
  \item Solving linear congruences.
\end{itemize}

\section*{Practice Problems}

\begin{itemize}
  \item Count or construct solutions to \(ax + by = c\).
  \item Normalize slope or direction pairs by gcd.
  \item Check whether two periodic processes ever align.
\end{itemize}

\section*{References}

\begin{itemize}
  \item \href{https://cp-algorithms.com/algebra/extended-euclid-algorithm.html}{cp-algorithms: Extended Euclidean Algorithm}.
  \item \href{https://usaco.guide/adv/extend-euclid?lang=cpp}{USACO Guide: Extended Euclid}.
  \item \href{https://codeforces.com/catalog}{Codeforces Catalog}.
\end{itemize}
